%\documentclass{cmspaper}
%\begin{document}
\section{Trigger Studies} \label{sec:trig}
% why I did it...
%All the trigger plots and explanations
%Quality of selection
%eff before and after analysis cuts, justification

The presence of two high energy electrons in the final state is exploited for the online selection 
of candidate leptoquark-antileptoquark events. 
Four electromagnetic High Level Trigger (HLT) paths are investigated
% \footnote{In the most recent HLT trigger table these names have been changed to $HLT\_EM80$, $HLT\_EM200$, 
%and $HLT\_IsoEle18\_L1R$.}
%In the most recent HLT trigger table these names have been chaged to HLT_EM80, HLT_EM200, and HLT_IsoEle18_L1R}
:
\begin{itemize}
\item the EM80 trigger 
\item the EM200 trigger 
\item the Photon15 trigger 
\item the Photon25 trigger .
\end{itemize}

The EM80 and EM200 triggers have been designed specifically to have high 
efficiency for electrons and photons of high transverse energy. 
They trigger on deposits of large electromagnetic transverse energy (supercluster) with
$E_T>80~$GeV and $E_T>200~$GeV respectively.

A combination of the EM80 or the EM200 trigger does result in a higher efficiency than the EM80 trigger 
alone due to the similar (identical?) isolation requirements of the two triggers.  Previous studies have shown
 this analysis may benefit from a trigger at higher transverse energy but looser isolation requirements.
 
 
The Photon 15 and Photon 25 triggers will be present in the early start up trigger menu, 
unlike the EM80 and EM200 triggers.  These trigger on any supercluster with 
$E_T>15~$GeV and $E_T>25~$GeV respectively with no requirements of track 
association nor isolation.  The higher level of background events passing this trigger is compensated 
for with the offline selection criteria.
%EMHighEt triggers on a high transverse energy electromagnetic deposit (supercluster) with
%$E_t>80~$GeV and loose isolation requirements (less than 5 GeV deposited in a 
%cone of $\Delta R<0.3$ \footnote{$\Delta R=\sqrt{\Delta\eta^2 + \Delta\phi ^2}$} surrounding the 
%primary supercluster in ECAL, less than 
%8 GeV in cone of $0.15<\Delta R<0.3$ in HCAL and less than 4 tracks in a cone $\Delta R<0.3$ of the tracker).
%EMVeryHighEt triggers on a very high transverse energy supercluster with $E_t>200~$GeV and no isolation 
%requirements. SingleElectronRelaxed selects an electron with $E_t>18~$GeV, an HCAL over ECAL energy ratio less than 5\%, 
%and $(1/E-1/p)<0.03$. More details on these triggers can be found in \cite{HeepHlt}.
% see https://twiki.cern.ch/twiki/bin/view/CMS/HEEPOnline
% https://twiki.cern.ch/twiki/bin/view/CMS/HighLevelTrigger06ExerciseDocumentation
% CMS AN 2007/009 and CMS Note 2006/078

The efficiencies for leptoquarks for these triggers have been studied with the FullSim 
samples of Section~\ref{sec:MCSamples} at two leptoquark mass points, and are summarized in Table~\ref{tab:HLTeffic} and 
Fig.~\ref{fig:HLTeffic}.

%At a luminosity of $10^{32}~cm^{-1}s^{-1}$ the estimated rates of the EMHighEt, EMVeryHighEt and SingleElectronRelaxed triggers 
%are 0.5, 0.1 and 10 Hz, respectively.
%In order to guarantee a high efficiency and at the same time contain the data sample size, the current decision is to perform 
%the analysis on the OR of the EMHighEt and EMVeryHighEt triggers.


  \begin{table}[htbp]
    \caption{\small \sl HLT efficiencies for LQ production at two mass values.}
    \label{tab:HLTeffic}
    \begin{center}
      \begin{tabular}{|llcc|} \hline
%               & \multicolumn{3}{c|}{Leptoquark Mass (GeV)} & 250 & 400 \\ \hline
        Leptoquark Mass (GeV)    & 250 & 400 \\ \hline
        EM80                                     & XXX$\pm$.012 & XXX$\pm$.008 \\
        EM200                                 & XXX$\pm$.013 & XXX$\pm$.015   \\ 
        PHOTON15                        & XXX$\pm$.012 & XXX$\pm$.013 &  \\ 
        PHOTON25                     & XXX$\pm$.011 & XXX$\pm$.006   \\ \hline
      \end{tabular}
    \end{center}
  \end{table}

\begin{figure}[htbp]
  \begin{center}
  \begin{tabular}{cc}
  \resizebox{12cm}{!}{\includegraphics{plots/UMD.eps}} 
  \end{tabular}
    \caption{\small \sl HLT efficiencies for LQ production at two mass values.}
    \label{fig:HLTeffic}
  \end{center}
\end{figure}




% add here the 1/30 of singleElecRelaxed?

%\end{document}

